\documentclass[11pt]{article}
% Кодировка и язык
\usepackage[utf8]{inputenc}           % Кодировка исходного файла (UTF-8)
\usepackage[T2A]{fontenc}             % Кодировка шрифтов для кириллицы
\usepackage[english,russian]{babel}   % Локализация: русские названия, переносы

% Поиск в PDF
\usepackage{cmap}                     

% Поля страницы
\usepackage{geometry}
\geometry{left=2cm, right=2cm, top=2cm, bottom=2cm}

% Математика
\usepackage{amsmath, amssymb, amsthm} % Базовые пакеты для математики
\usepackage{mathtools}                % Расширение amsmath (доп. символы и окружения)

% Картинки
\usepackage{graphicx}                 % Вставка картинок (\includegraphics)
\graphicspath{{pictures/}}            % Путь к папке с картинками
\usepackage{float}                    % Позволяет использовать [H] для жёсткого позиционирования
\usepackage{wrapfig}                  % Обтекание текста вокруг картинок

% Таблицы
\usepackage{array}                    % Расширенные возможности для таблиц
\usepackage{booktabs}                 % Красивые таблицы (\toprule, \midrule, \bottomrule)
\usepackage{multirow}                 % Объединение строк в таблицах
\usepackage{tabularx}                 % Таблицы с автоматической шириной колонок
\usepackage{ragged2e}                 % Улучшенное выравнивание (используется в tabularx)

% Списки
\usepackage{enumitem}                 % Гибкая настройка списков (itemize, enumerate)

% Колонтитулы
\usepackage{fancyhdr}                 

% Гиперссылки
\usepackage{xcolor}                   % Работа с цветами
\usepackage{hyperref}                 % Гиперссылки, оглавление, перекрёстные ссылки

% Умные ссылки
\usepackage{cleveref}                 

% Типографика
\usepackage{microtype}                % Улучшает переносы и межсловные расстояния

% Подписи к картинкам и таблицам
\usepackage{caption}                  % Настройка подписей (шрифт, отступы)

% Рисование схем и графиков
\usepackage{tikz}                     % Рисование схем, графиков, диаграмм
\usepackage{pgfplots}                 % Построение графиков функций (на основе tikz)
\pgfplotsset{compat=newest}           % Версия совместимости pgfplots

% Алгоритмы и код
\usepackage{algorithm}                % Окружения для алгоритмов
\usepackage{algorithmic}              % Синтаксис для алгоритмов
\usepackage{listings}                 % Вставка кода с подсветкой синтаксиса


% Настройка нумерации секций и подсекций
\renewcommand{\thesection}{\arabic{section}}                   
\renewcommand{\thesubsection}{\thesection.\arabic{subsection}} 


% Настройка списков: убираем лишние отступы слева
\setlist[itemize]{leftmargin=}
\setlist[enumerate]{leftmargin=}


% Определение окружений для теорем, лемм, определений и т.д.
% Нумерация привязана к номеру секции
\newtheorem{theorem}{Теорема}[section]
\newtheorem{definition}{Определение}[section]
\newtheorem{remark}{Замечание}[section]
\newtheorem{lemma}{Лемма}[section]
\newtheorem{example}{Пример}[section]
\newtheorem{consequence}{Следствие}[section]
\newtheorem{algorithmm}{Алгоритм}
\newtheorem{question}{Вопрос}[section]
\newtheorem{answer}{Ответ}[section]
\newtheorem{statement}{Утверждение}[section]


% Определение цветов
\definecolor{pblue}{rgb}{0.13,0.13,1}    
\definecolor{pgreen}{rgb}{0,0.5,0}       
\definecolor{pred}{rgb}{0.9,0,0}         
\definecolor{pgrey}{rgb}{0.46,0.45,0.48}

\newcommand{\lectureNumber}{4}
\newcommand{\lectureTitle}{Линейные ДУ $n$-ого порядка с постоянными коэффициентами в случае кратных корней. ФСР линейного однородного ДУ $n$-ого порядка с постоянными коэффициентами. Линейные неоднородные ДУ $n$-ого порядка с постоянными коэффициентами. Квазимногочлены.}
\newcommand{\courseName}{Элементы дифференциальных уравнений}
\newcommand{\semester}{3 семестр, 2026/2027 уч. год}

\pagestyle{fancy}
\fancyhf{}
\fancyhead[L]{Тот факт, что будильник не прозвенел, изменил уже много человеческих судеб}
\fancyhead[R]{\thepage}
\setlength{\headheight}{15pt}

\hypersetup{
    colorlinks=true,
    linkcolor=blue!70!black,
    citecolor=blue!70!black,
    filecolor=blue!70!black,
    urlcolor=blue!70!black,
}

\begin{document}


\begin{titlepage}
    \centering
    
    {\large Московский физико-технический институт \par}
    {\large Высшая школа программной инженерии \par}
    \vspace{2cm}

    {\Large \textbf{\courseName} \par}
    {\large \semester \par}
    
    \vspace*{\fill}

    {\Huge \textbf{Лекция \lectureNumber} \par}
    {\Huge \textbf{\lectureTitle} \par}
    
    \vspace*{\fill}
    
    {\large Лектор: Саулин Сергей Михайлович \par}
    {\large Автор: \href{https://t.me/ceagest}{Юра} \par}

\end{titlepage}

\pagenumbering{arabic}

\section{Веселье продолжается}

В прошлой лекции мы рассматривали наши любимые линейные ДУ $n$-ого порядка с постоянными коэффициентами в случае кратных корней. Вот, кстати, это самое уравненьице:
\[a_0 \cdot y^{(n)} + a_1 \cdot y^{(n - 1)} + \dots + a_{n - 1} \cdot y' + a_n \cdot y = 0, \text{ где } a_0 \neq 0, \ a_1, a_2, \dots, a_n \in \mathbb{C} \tag{1} \label{eq:1}\]
С ним рука об руку идёт характеристический многочлен:
\[\chi(\lambda) = a_0 \cdot \lambda^n + a_1 \cdot \lambda^{n-1} + \dots + a_n \tag{2} \label{eq:2}\]
В доказательстве остановились на самом интересном месте, но, чтобы не утруждать читателя переходом в соседний затех, приведу ту часть доказательства и здесь, а потом незаметно продолжу писать его дальше :)

\begin{theorem}
    Пусть $\lambda_1, \lambda_2, \dots, \lambda_m$ --- попарно различные корни характеристического уравнения $\chi(\lambda) = 0$, алгебраические
    кратности которых равны $k_1, k_2, \dots, k_m$ соответственно. $\displaystyle \left(\sum_{i = 1}^{m} k_i = n\right)$. Тогда общее решение
    уравнения (\ref{eq:1}) имеет следующий вид: \[\tag{3} \label{eq:3} y(x) = \sum_{j = 1}^{m} P_j(x) \cdot e^{\lambda_j x}, \text{ где } P_j(x) 
    \text{ --- многочлен степени } \leq k_j - 1, \ j = 1, \dots, m\]
\end{theorem}

\begin{proof}
    Покажем, что любая функция вида (\ref{eq:3}) является решением уравнения (\ref{eq:1}). Уравнение (\ref{eq:1}) можно записать в следующем виде:
    \[\chi(\mathcal{D})(y(x)) = 0\] Достаточно убедиться, что решениями являются следующие функции: \[e^{\mu x}, x \cdot e^{\mu x}, \dots, x^{k-1}
    \cdot e^{\mu x}, \quad \text{где } \mu \text{ --- корень } \chi(\lambda) = 0 \text{ кратности } k\]
    Тогда \[\chi(\lambda) = \zeta(\lambda) \cdot (\lambda - \mu)^k, \quad  \text{где } \zeta(\lambda) \text{ --- многочлен, } \zeta(\mu) \neq 0
    \text{ (иначе иная кратность у } \mu)\] Следовательно, по формуле сдвига имеем \[\chi(\mathcal{D})(e^{\mu x} \cdot x^j) = e^{\mu x} 
    \chi(\mathcal{D} + \mu \cdot I)(x^j) = e^{\mu x} \zeta(\mathcal{D + \mu \cdot I}) \cdot \mathcal{D}^k(x^j) = 0 \quad \forall j < k\]
    То есть, любая функция вида (\ref{eq:3}) является решением уравнения (\ref{eq:1}). \\ 
    Теперь докажем, что любое решение уравнения (\ref{eq:1}) имеет вид (\ref{eq:3}): определим следующие функции для каждого $\lambda_j$
    \[y_1(x) = e^{\lambda_1 x}, \ y_2(x) = x \cdot e^{\lambda_1 x}, \dots, \ y_{k_1} (x) = x^{k_1 - 1} e^{\lambda_1 x}\]
    \[y_{k_1 + 1}(x) = e^{\lambda_2 x},\  y_{k_1 + 2}(x) = x \cdot e^{\lambda_2 x}, \dots, \ y_{k_1 + k_2}(x) = x^{k_2 - 1} \cdot e^{\lambda_2 x}\]
    \[\vdots\] \[y_{k_1+...+k_{m-1}+1}(x) = e^{\lambda_m x}, \ y_{k_1+...+k_{m-1}+2}(x) = x \cdot e^{\lambda_m x}, \dots, 
    \ y_{k_1+...+k_m}(x) = y_n(x) = x^{k_m - 1} \cdot e^{\lambda_m x}\]
    Пусть $y(x)$ --- решение уравнения (\ref{eq:1}) с начальными условием \[y(x_0) = y_0, \ y'(x_0) = y_0', \dots, \ y^{(n-1)}(x_0) = y_0^{(n-1)}\]
    Покажем, что \[\exists c_1, c_2, \dots, c_n \in \mathbb{C}: \ y(x) = \sum_{k = 1}^{n} c_k \cdot y_k(x)\] 
    Многократно дифференцируя равенство выше и подставляя точку $x_0$, получаем следующую системку (не забываем про обозначения из начальных условий выше):
    \[y_0^{(j)} = \sum_{k = 1}^{n} c_k \cdot y_k^{(j)}(x_0), \quad j = 0, 1, \dots, n-1 \tag{4} \label{eq:4}\]
    Рассмотрим же матрицу системы (\ref{eq:4}): \[M = \begin{pmatrix}
        y_1(x_0) & \dots & y_n(x_0) \\ y_1'(x_0) & \dots & y_n'(x_0) \\ \vdots & \ddots & \vdots \\ y_1^{(n-1)}(x_0) & \dots & y_n^{(n-1)} (x_0)
    \end{pmatrix}\]
    Если эта матрица невырожденна, то всё доказано --- обратив её, получаем искомые коэффициенты $c_k$. Пусть $M$ вырожденная, тогда ее строки линейно зависимы, следовательно 
    \[\tag{5} \label{eq:5} \exists b_0, b_1, \dots, b_{n-1} \in \mathbb{C}: \ \sum_{i = 0}^{n - 1} y_j^{(i)} \cdot b_{n-1-i} = 0 \quad \forall j = 1, \dots, n\]
    Рассмотрим многочлен \[P(\lambda) = \sum_{i = 0}^{n-1} \lambda^i \cdot b_{n-1-i}\]
    Понятно, что его степень не превосходит $n - 1$ (если коэффициент при старшей степени не ноль, то ровно $n - 1$, иначе строго меньше). Следовательно, по основной теореме алгебры количество его корней также не превосходит $n - 1$. \\
    Пользуясь нашим вспомогательным многочленом, перепишем (\ref{eq:5}) в эквивалентной операторной форме:
    \[\tag{6} \label{eq:6} P(\mathcal{D})(y_j) \bigg|_{x = x_0} = 0 \quad \forall j = 1, \ldots, n\]
    Теперь, вспомним чудесную леммочку, доказанную в прошлой серии:
    \begin{lemma}
        Пусть $x_0 \in \mathbb{R}$ фиксированно, $\mu \in \mathbb{C}$, а также \[P(\mathcal{D}) \left(e^{\mu x} x^k\right) \bigg|_{x=x_0} = 0 
        \quad \forall k \in \{0, 1, \dots, s\}\]
        Тогда $\mu$ --- корень многочлена $P$ кратности $\geq s + 1$.
    \end{lemma} \leavevmode
    \\
    Заметим, что $y_j$ --- это функции в точности нужного нам вида. Если мы запишем условие (\ref{eq:6}) для $j = 1, \ldots, k_1$, то из леммочки немедленно получим, что $\lambda_1$ --- корень многочлена $P$ кратности $\geq k_1$. Продолжая так делать и дальше, получим, что:
    \[\lambda_2 \text{ имеет кратность } \geq k_2\]
    \[\vdots\]
    \[\lambda_m \text{ имеет кратность } \geq k_m\]
    Следовательно, у многочлена $P$ корней не меньше, чем $k_1 + k_2 + \ldots + k_m = n$. Но $\deg P \leq n - 1$, откуда получаем противоречие. \\
    Таким образом, мы доказали, что справедливо равенство (\ref{eq:4}), то есть начальные условия линейным образом выражаются через функции $y_j$, посчитанные в точке $x_0$. Следовательно, по теореме о существовании и единственности решения ЗК, функции $y(x)$ и $\sum_{k = 1}^n c_k \cdot y_k(x)$ совпадают не только в точке $x_0$, но и в целой её окрестности. Тогда по теореме о продолжении решения эти функции совпадают вообще на всей числовой прямой, что и требовалось.
\end{proof}

\section{Овеществление комплексных решений}
Теперь встаёт закономерный вопрос об овеществлении решений. Ответит нам на него одноимённая теорема, доказательство которой почти полностью аналогично тому, которое мы проделывали для случая попарно различных корней.
Здесь мы также предполагаем, что все коэффициенты $a_i$ уравнения (\ref{eq:1}) \textbf{вещественны}.

\begin{theorem}
    Пусть $\lambda_1, \ldots, \lambda_m$ --- попарно различные корни характеристического уравнения, алгебраические кратности которых равны $k_1, \ldots, k_m$ соответственно ($k_1 + \ldots + k_m = n$). Предположим, что 
    \[\lambda_{2j - 1} = \alpha_j + i\beta_j, \; \; \lambda_{2j} = \overline{\lambda_{2j - 1}} \quad \forall j = 1, \ldots, l\]
    \[\lambda_{2l + 1}, \ldots, \lambda_m \in \mathbb{R}\]
    Тогда общее \textbf{вещественное} решение уравнения (\ref{eq:1}) имеет следующий вид:
    \[y(x) = \sum_{j = 1}^l e^{\alpha_j x} \cdot (P_j(x) \cos \beta_j x + Q_j(x) \sin \beta_j x) + \sum_{s = 2l + 1}^m R_s(x) e^{\lambda_s x}\]
    Здесь $P_j$ и $Q_j$ --- многочлены с вещественными коэффициентами, степени которых $\leq k_j - 1$ при $j = 1, \ldots, l$, а $R_s$ --- многочлен степени $\leq k_s - 1$ при $s = (2l + 1), \ldots, m$.
\end{theorem}

\section{Фундаментальная система решений линейного однородного ДУ $n$-ого порядка с постоянными коэффициентами}

\begin{definition}
    Говорят, что набор функций $z_1, z_2, \ldots, z_d \colon \mathbb{R} \to \mathbb{C}$ --- фундаментальная система решений уравнения (\ref{eq:1}), если
    \begin{enumerate}
        \item $z_1(x), z_2(x), \ldots, z_d(x)$ --- линейно независимы, то есть
        \[\sum_{k = 1}^d c_k z_k(x) = 0 \; \; \forall x \in \mathbb{R} \implies c_1 = c_2 = \ldots = c_d = 0\]
        \item Любое решение уравнения (\ref{eq:1}) является линейной комбинацией $z_1(x), z_2(x), \ldots, z_d(x)$.
    \end{enumerate}
\end{definition} \leavevmode
\\
Доказывая теорему $1.1$, мы на самом деле показали, что пространство решений (\ref{eq:1}) конечномерно, и базисом в нём можно взять как раз функции $y_j$ (легко показать, что они линейно независимы, так как они суть экспонента на моном). Таким образом, мы установили

\begin{consequence}
    ФСР уравнения (\ref{eq:1}) существует. Более того, она состоит из $n$ функций (то есть $d = n$).
\end{consequence}

\section{Линейные неоднородные ДУ $n$-ого порядка с постоянными коэффициентами}

Теперь будем рассматривать уравнения следующего вида:

\[\tag{7} \label{eq:7} a_0 y^{(n)} + a_1 y^{(n - 1)} + \ldots + a_n y = F(x), \; \; a_0, a_1, \ldots, a_n \in \mathbb{C}, \; \; a_0 \neq 0\]
\[F \colon I \to \mathbb{R}, \; \; I \subset \mathbb{R} \text{ --- интервал}\]

\begin{remark}
    Следующие два утверждения позволяют чуть больше понять, как же устроено общее решение (\ref{eq:7}):
    \begin{enumerate}
        \item Пусть $y_1(x)$ --- решение \textbf{соответствующего} однородного уравнения (\ref{eq:1}), а $y_2(x)$ --- решение уравнения (\ref{eq:7}). Тогда $y(x) = y_1(x) + y_2(x)$ --- решение уравнения (\ref{eq:7}).
        \item Пусть $y_1(x), y_2(x)$ --- решения уравнения (\ref{eq:7}). Тогда $y_1(x) - y_2(x)$ --- решение уравнения (\ref{eq:1}).
    \end{enumerate}
\end{remark}

\begin{proof}
    Докажем утверждения по очереди.
    \begin{enumerate}
        \item Так как $y_1(x)$ --- решение (\ref{eq:1}), то 
        \[a_0 y_1^{(n)} + \ldots + a_n y_1 = 0\]
        Так как $y_2(x)$ --- решение (\ref{eq:7}), то 
        \[a_0 y_2^{(n)} + \ldots + a_n y_2 = F(x)\]
        Следовательно,
        \[a_0 (y_1 + y_2)^{(n)} + \ldots + a_n (y_1 + y_2) = F(x)\]
        \item Так как $y_1(x)$ и $y_2(x)$ --- решения (\ref{eq:7}), то 
        \[a_0 y_1^{(n)} + \ldots + a_n y_1 = F(x)\]
        \[a_0 y_2^{(n)} + \ldots + a_n y_2 = F(x)\]
        Следовательно,
        \[a_0 (y_1 - y_2)^{(n)} + \ldots + a_n (y_1 - y_2) = 0\]
    \end{enumerate}
\end{proof} \leavevmode
\\
Из этого замечания сразу следует, что общее решение неоднородного уравнения (\ref{eq:7}) есть сумма частного решения (\ref{eq:7}) и общего решения однородного уравнения (\ref{eq:1}). Таким образом, задача об изучении множества решений (\ref{eq:7}) свелась к задаче о поиске частного решения (\ref{eq:7}). Этим мы и будем заниматься дальше.

\subsection{Квазимногочлены}

Будем обозначать за $\mathbb{R}_n[x]$ пространство многочленов с вещественными коэффициентами от переменной $x$ степени $\leq n$. Также условимся, что $\mathbb{R}[x]$ --- кольцо многочленов над полем $\mathbb{R}$, то есть
\[\mathbb{R}[x] = \bigcup_{n \geq 0} \mathbb{R}_n[x]\]

\begin{definition}
    \textbf{Квазимногочленом} $F$ степени $m$ с показателем $\lambda = \alpha + i \beta$ называется функция вида
    \[F(x) = e^{\alpha x} (p(x) \cos \beta x + q(x) \sin \beta x), \; \; p, q \in \mathbb{R}[x], \; \; \max \{\deg p, \deg q\} = m\]
\end{definition}

\begin{definition}
    Определим \textbf{пространство квазимногочленов} фиксированной степени и фиксированного показателя:
    \[Q_{\lambda, m} = \{r(x) = e^{\alpha x} (p(x) \cos \beta x + q(x) \sin \beta x) \; | \; p, q \in \mathbb{R}_{m - 1}[x]\} \]
\end{definition} \leavevmode
\\
Из определения вытекает очевидное наблюдение:
\[Q_{\overline{\lambda}, m} = Q_{\lambda, m}\]

\begin{definition}
    Для дальнейшей работы нам понадобится ещё и такое пространство:
    \[Q_{\lambda, m}^k = \{x^k r(x) \; | \; r \in Q_{\lambda, m}\}, \; \; k \in \mathbb{Z}_{+}\]
\end{definition}

\begin{remark}
    Отметим следующие милые тривиальные свойства:
    \begin{enumerate}
        \item \[Q_{\lambda, m}^k \subset Q_{\lambda, m + k}\]
        \item \[Q_{\overline{\lambda}, m}^k = Q_{\lambda, m}^k\]
    \end{enumerate}
\end{remark}

\begin{remark}
    Если $\lambda = \alpha$, то есть $\beta = 0$, то
    \[Q_{\lambda, m} = \left\langle e^{\alpha x} x^t \; | \; t = 0, 1, \ldots, (m - 1)\right\rangle \]
    Если же $\lambda = \alpha + i \beta$, то есть $\beta \neq 0$, то 
    \[Q_{\lambda, m} = \left\langle e^{\alpha x} x^t \cos \beta x, e^{\alpha x} x^s \sin \beta x \; | \; t, s = 0, \ldots, (m - 1)\right\rangle \]
\end{remark}

\begin{exercise}
    Покажите, что все функции, входящие в набор
    \[\{e^{\alpha x} x^t \cos \beta x, e^{\alpha x} x^s \sin \beta x \; | \; t, s = 0, \ldots, (m - 1)\}\]
    являются линейно независимыми.
\end{exercise}

\begin{theorem}
    Пусть $F \in Q_{\mu, m}$, $\mu \in \mathbb{C}$, а $s$ --- кратность $\lambda = \mu$ как корня характеристического уравнения $\chi(\lambda) = 0$ ($s = 0$, если $\lambda = \mu$ не является корнем). Тогда у уравнения (\ref{eq:7}) есть частное решение, которое принадлежит $Q_{\mu, m}^s$.
\end{theorem}

\begin{proof}
    Проведем доказательство для $\mu \in \mathbb{R}$. Пусть $\mu$ --- корень $\chi(\lambda) = 0$ кратности $s \geq 0$. Тогда
    \[\chi(\lambda) = \zeta(\lambda)(\lambda - \mu)^s\]
    Здесь многочлен $\zeta(\mu) \neq 0$. Тогда уравнение (\ref{eq:7}) эквивалентно
    \[\tag{8} \label{eq:8} \chi(\mathcal{D})(y(x)) = F(x)\]
    Но
    \[\chi(\mathcal{D}) = \zeta(\mathcal{D})(\mathcal{D} - \mu \mathcal{I})^s\]
    Покажем, что отображение $\chi(\mathcal{D}) \colon Q_{\mu, m}^s \to Q_{\mu, m}$ является биекцией. Тогда, обратив его, при помощи равенства (\ref{eq:8}) получим искомое частное решение. \\
    Достаточно показать, что каждое из отображений 
    \[(\mathcal{D} - \mu \mathcal{I})^s \colon Q_{\mu, m}^s \to Q_{\mu, m}\]
    \[\zeta(\mathcal{D}) \colon Q_{\mu, m} \to Q_{\mu, m}\]
    является биекцией в силу представления выше. \\
    \textbf{To be continued...}
\end{proof}

\end{document}